\documentclass[10pt,a4paper]{amsart}
\usepackage[margin=19mm]{geometry}
\usepackage{amsmath,amssymb,mathtools}
\usepackage[scr=rsfs,cal=euler]{mathalfa}
\usepackage{xfrac}
\usepackage[unicode,hidelinks,bookmarks=false]{hyperref}
\newcommand{\ie}{i.e.\ }
\newcommand{\cf}{cf.\ }
\newcommand{\resp}{respectively\ }
\newcommand{\Subsection}{\subsection}
\providecommand{\nobreakdash}{\nobreak\hbox{-}}
\setlength{\emergencystretch}{2em}
\multlinegap0pt
\usepackage{bm}
\usepackage{bbm}
\usepackage{dsfont}
\usepackage{xspace}
\usepackage{cancel}
\usepackage{mathtools}
\usepackage{amsthm}
\makeatletter
\@ifundefined{theorem}{\newtheorem{theorem}{Theorem}}{}
\makeatother
\title[Cop number of partial cubes]{Cop number of partial cubes}
\author{Nicholas Crawford, Vesna Ir\v{s}i\v{c} Chenoweth}
\date{Source: arXiv:2510.06832v1 (CC BY 4.0)}
\begin{document}
\maketitle

\noindent\emph{Source and licence.} Cop number of partial cubes. Version arXiv:2510.06832v1 (CC BY 4.0), distributed under the Creative Commons Attribution 4.0 International licence, as recorded in the arXiv licence metadata for this exact version and shown on the abstract page https://arxiv.org/abs/2510.06832v1. Full rights notice: licenses/arxiv-2510.06832-CC-BY-4.0.txt.\par\smallskip

\noindent\textit{Editorial scope.} Counted unit: the Theorem whose source label is \texttt{thm:n/3} in the article (source0.tex lines 358--361, source label \texttt{thm:n/3}), delivered together with the article's own complete original proof of it (source0.tex lines 363--388, one \texttt{proof} environment of 26 lines). No mathematics is written, repaired or re-derived: statement and proof are the article's own text line for line. Required context retained uncounted: none. Every definition, display and label of those lines is retained unchanged. Omitted from this package and not redistributed: 1--357, 362--362, 389--544. Nothing inside a retained excerpt is omitted or altered: every retained body is delivered exactly as the article's lines are written, so no omission receipt of any kind accompanies this package. Citations: the retained excerpts carry no citation command at all, so no bibliography entry of the article is transcribed and no reference is redistributed. Reference adaptation: none was needed and none was made; every reference inside the delivered unit resolves inside it, so no unresolved reference or citation is produced. Printed slips kept literal: no printed word, symbol, hypothesis or number of the counted statement or of its proof was altered. Layout and class: the article's own class is \texttt{article} with packages including babel, amsmath, amssymb, amsfonts, amsthm, tikz, color, pgf, graphicx, url, enumerate, xcolor, hyperref, natbib; the wrapper is the lane's pinned \texttt{amsart} editorial wrapper at 10pt on A4, which restates the theorem-like environments the retained text needs and adds the article's own macro definitions that the retained text needs (none), with extra wrapper packages (bm, bbm, dsfont, xspace, cancel, mathtools). The delivered numbering is the unit's own. Assets: no retained span contains a figure environment, a graphics-inclusion command, a PDF-page inclusion, a table or a TikZ picture; no image file is copied into this package, the package declares no asset dependency and no image file is redistributed with it. Licence: version arXiv:2510.06832v1 of this article is distributed under the Creative Commons Attribution 4.0 International licence, as recorded in the arXiv licence metadata for this exact version and shown on the abstract page https://arxiv.org/abs/2510.06832v1. A full rights notice is delivered at licenses/arxiv-2510.06832-CC-BY-4.0.txt. Characters: every retained excerpt is pure ASCII, so the delivered engine, XeLaTeX with its default Latin Modern text font, reports no missing character.
\begin{theorem}
    \label{thm:n/3}
    If $n \geq 6$, then $c(\Gamma_n)\leq \lceil \frac{n}{3}\rceil$.
\end{theorem}

\begin{proof}
    Let $k=\lceil \frac{n}{3}\rceil$. We provide a winning strategy for $k$ cops $c_1, \ldots, c_k$. Partition each binary string of length $n$ into $k$ blocks $B_1, \ldots, B_k$ of length 3 except $B_k$ which is of length $n - 3(k-1) \in \{1,2,3\}$. More precisely, the blocks of the string $x_1 \ldots x_n$ are $B_i = x_{3i-2} x_{3i-1} x_{3i}$ for $i \in [k-1]$ and $B_k \in \{x_n, x_{n-1} x_n, x_{n-2} x_{n-1} x_n\}$. In our strategy, which consists of different phases, cop $c_i$ is associated with block $B_i$, and all cops start in $0^n$. If cop $c_i$ and the robber have the same value on $B_j$ we say that $c_i$ matches the robber on $B_j$.

    Before presenting the strategy, observe that vertices with all possible values of $B_i$ where values on other blocks are fixed induce a subgraph of $\Gamma_n$ that is isomorphic to $\Gamma_3$ or an induced subgraph of $\Gamma_3$, and that $000$ is at distance at most 2 from all other possible values of $B_i$ and at distance exactly 2 only from $101$. When describing the cop's strategy to match the robber on some $B_i$ we will say that the cop moves closer to the robber on $B_i$, which means that the cop moves closer to the robber in the corresponding $\Gamma_3$. Saying that a player moves on $B_i$ will mean that they change a bit from $B_i$.

    \begin{description}
        \item[Phase 1.] Until it holds for every $i \in [k]$ except maybe one that $c_i$ matches the robber on $B_i$.\\
        The strategy for cop $c_i$ during Phase 1 is to only change bits in $B_i$, aiming to match the robber there (the remaining bits stay $0$). Once they match, the cop can maintain this property by moving on $B_i$ if and only if the robber moves on $B_i$, and not moving otherwise. If the robber moves on $B_j$, $j \neq i$, then the cop $c_i$ moves closer to the robber on $B_i$ (decreasing the distance between them by 1). If the robber moves on $B_i$, $c_i$ moves such that the distance between the cop and the robber does not increase. If the robber does not move on $B_i$ for at least two rounds, $c_i$ can match the robber on $B_i$. Since $k \geq 3$, there is at most one block $B_i$ on which $c_i$ is not able to match the robber. Without loss of generality, let this be the block $B_k$.

        \item[Phase 2.] For every $i \in [k]$, the cop $c_i$ matches the robber on $B_i$.
        \begin{description}
            \item[Phase 2.1.] The game never reaches the condition from Phase 2.\\
            By the strategy from Phase 1, this is only possible if the robber always moves on $B_k$ (except maybe once). Thus, cops $c_1, \ldots, c_{k-1}$ can match the robber on $B_1$ by moving closer to the robber there. Then they can also match the robber on $B_2$, etc., until they all match the robber on blocks $B_1, \ldots, B_{k-1}$. Note that by matching the robber on blocks sequentially, we avoid the possible problem of the cop wanting to move to a vertex outside of $V(\Gamma_n)$. As $k \geq 3$, we thus have at least two cops that match the robber everywhere except on $B_k$. The final part of their strategy is to imagine the game on $\Gamma_{|B_k|}$ according to the robber's moves on $B_k$ and as $c(\Gamma_m) \leq 2$ for $m \in [3]$, these two cops can win the game. Once they do, the robber is caught on $\Gamma_n$ too.

            Note that during this phase, cop $c_i$ always matches the robber on $B_i$ for every $i \in [k-1]$. Observe also that if on some $B_j$ cop $c_i$ is at distance 2 from the robber anytime during Phase 2.1, then $c_i$ is on $000$ and the robber is on $101$. 

            \item[Phase 2.2.] After the condition from Phase 2 is established.\\
            Suppose the robber moves on $B_j$. If the cop $c_i$ matches the robber on $B_j$, then $c_i$ simply moves on $B_j$ as well, so that he still matches the robber. If $c_i$ does not yet match the robber on $B_j$, then we distinguish between the following cases.
            \begin{itemize}
                \item If robber is at distance 1 from $c_i$ on $B_j$, $c_i$ moves such that he matches the robber on $B_j$ after this move.
                \item If robber is at distance $2$ from $c_i$ on $B_j$, then as observed in Phase 2.1., $c_i$ is on $000$ and robber is on $101$. Cop's strategy is to move to $001$. As on $B_{j+1}$, $c_i$ either matches the robber (who is on $101$ on $B_j$) or is on $000$, this is a legal move in $V(\Gamma_n)$. If in the next round, the robber moves on $B_j$ again and does not match $c_i$, then $c_i$ returns to $000$. (To avoid forcing the cop to move outside of $V(\Gamma_n)$ when trying to match the robber on $B_{j+1}$.) Otherwise, if the robber moves on some $B_\ell$, $\ell \neq j$, then $c_i$ can match the robber on $B_j$ after his move. 
            \end{itemize}
            This strategy ensures that if the robber moves on some $B_j$ and then moves on a different block in the next round (or does not move at all), all cops match the robber on $B_j$ after these two rounds. So, unless the robber is moving only on one block for the rest of the game, he is caught in a finite number of rounds. If the robber keeps moving on only one block, say on $B_\ell$, then first wait for all cops except $c_\ell$ to match the robber on blocks $B_i$, $i \in [k] \setminus \{\ell\}$. Afterwards, as $\Gamma_3$ has only one vertex at a distance 2 from $000$ and the robber is always moving on $B_\ell$, either in this or the next round, the robber is on $001$ or $100$, so all cops but $c_\ell$ now also match the robber on $B_\ell$, thus, the robber is caught. \hfill \qedhere
        \end{description}
    \end{description}
\end{proof}

\end{document}
